\section{Results}

All numbers below come from the saved runs in \code{eval/results/} and are collected by
\code{eval/summarize.py} into \code{eval/RESULTS\_v2.md}. Haiku results are the mean $\pm$ standard
deviation over three runs; Qwen decodes greedily, so it has one run. The 95\% intervals come from a
bootstrap over the questions (\cref{sec:grading}).

\subsection{E1: accuracy on our question set}

\Cref{tab:e1-v2} and \cref{fig:e1-v2} show the accuracy on the 50 test questions. ``Wrong decl.''
counts questions the system refused although it should have answered (out of 42), ``Wrong'' the
answers that are simply wrong.

\begin{table}[ht]
\centering\small
\setlength{\tabcolsep}{3.5pt}
\resizebox{\linewidth}{!}{%
\begin{tabular}{llrrrrrr}
\toprule
System & Model & Accuracy \% & 95\% CI & Wrong decl. & Wrong & Tokens in / out & \$ / 1k q \\
\midrule
A workflow       & Haiku 5.5   & \textbf{96.0 $\pm$ 0.0} & 90--100 & 0.0 & 2.0 & 2,803 / 305 & 0.43 \\
A0 template      & Haiku 5.5   & 96.0 $\pm$ 0.0          & 90--100 & 0.3 & 1.7 & 2,384 / 204 & 0.34 \\
A2 workflow v2   & Haiku 5.5   & 96.0 $\pm$ 0.0          & 90--100 & 0.0 & 2.0 & 3,897 / 305 & 0.54 \\
C simple RAG     & Haiku 5.5   & 68.7 $\pm$ 1.2          & 56--82  & 0.3 & 15.3 & 3,583 / 462 & 0.59 \\
N closed book    & Haiku 5.5   & 30.7 $\pm$ 1.2          & 20--43  & 11.7 & 23.0 & 112 / 796  & 0.41 \\
N closed book    & Opus 5.5    & 62                      & 48--74  & 7   & 12   & 112 / 819  & 16.83 \\
\midrule
A workflow       & Qwen2.5-7B  & \textbf{84}  & 74--94 & 5  & 3  & 1,330 / 112 & 0 \\
A0 template      & Qwen2.5-7B  & 84           & 74--94 & 5  & 3  & 1,049 / 77  & 0 \\
C simple RAG     & Qwen2.5-7B  & 38           & 24--52 & 3  & 28 & 2,472 / 137 & 0 \\
N closed book    & Qwen2.5-7B  & 14           & 6--24  & 33 & 10 & 88 / 263    & 0 \\
\midrule
A workflow       & Qwen2.5-3B  & \textbf{62}  & 48--76 & 8  & 11 & 1,394 / 126 & 0 \\
A0 template      & Qwen2.5-3B  & 62           & 48--76 & 8  & 11 & 1,124 / 98  & 0 \\
C simple RAG     & Qwen2.5-3B  & 40           & 26--54 & 2  & 28 & 2,472 / 194 & 0 \\
N closed book    & Qwen2.5-3B  & 18           & 8--30  & 34 & 7  & 88 / 127    & 0 \\
\midrule
A workflow       & Qwen2.5-3B (EC2 CPU) & 54  & 40--66 & 14 & 9 & 1,052 / 131 & 0 \\
\bottomrule
\end{tabular}}
\caption{E1 on our question set (50 test questions). Haiku: mean $\pm$ standard deviation over three
runs (wrong declines and wrong answers are means); Qwen: one run. Tokens are the mean per question;
the last column is the cost per 1,000 questions. Local models cost nothing per question. The last
row runs Qwen2.5-3B with Ollama on the CPU of the EC2 instance.}
\label{tab:e1-v2}
\end{table}

\begin{figure}[ht]
\centering
\begin{tikzpicture}
\begin{axis}[
  ybar, bar width=9pt, width=0.85\linewidth, height=55mm,
  ymin=0, ymax=110, ytick={0,20,40,60,80,100}, ylabel={Accuracy (\%)},
  symbolic x coords={Haiku 5.5,Qwen2.5-7B,Qwen2.5-3B}, xtick=data,
  enlarge x limits=0.25, legend style={at={(0.5,1.03)},anchor=south,legend columns=-1,font=\small,/tikz/every even column/.append style={column sep=3mm}},
  nodes near coords, every node near coord/.append style={/pgf/number format/precision=0, font=\scriptsize\bfseries, color=white, anchor=north, yshift=-1pt},
  error bars/y dir=both, error bars/y explicit, error bars/error bar style={black!60}]
\addplot[fill=wf,draw=none] coordinates {(Haiku 5.5,96) -= (0,6) += (0,4) (Qwen2.5-7B,84) -= (0,10) += (0,10) (Qwen2.5-3B,62) -= (0,14) += (0,14)};
\addplot[fill=rag,draw=none] coordinates {(Haiku 5.5,68.7) -= (0,12.7) += (0,13.3) (Qwen2.5-7B,38) -= (0,14) += (0,14) (Qwen2.5-3B,40) -= (0,14) += (0,14)};
\addplot[fill=cb,draw=none] coordinates {(Haiku 5.5,30.7) -= (0,10.7) += (0,12.3) (Qwen2.5-7B,14) -= (0,8) += (0,10) (Qwen2.5-3B,18) -= (0,10) += (0,12)};
\legend{A workflow, C simple RAG, N closed book}
\end{axis}
\end{tikzpicture}
\caption{E1 accuracy on our question set by system and model. The error bars are the 95\% bootstrap
intervals over the 50 questions.}
\label{fig:e1-v2}
\end{figure}

\Cref{tab:e1-cat} breaks the results down by category.

\begin{table}[ht]
\centering\small
\begin{tabular}{llrrrrrr}
\toprule
System & Model & Lookup & Change & Paraphrase & Why & Decline & Latest \\
 & & (10) & (10) & (8) & (12) & (8) & (2) \\
\midrule
A & Haiku 5.5  & 10 & 10 & 8 & 11 & 7 & 2 \\
C & Haiku 5.5  & 6  & 6  & 3 & 10.3 & 8 & 1 \\
N & Haiku 5.5  & 2.3 & 1.7 & 0 & 5.7 & 5.7 & 0 \\
\midrule
A & Qwen2.5-7B & 8  & 10 & 8 & 7  & 7 & 2 \\
C & Qwen2.5-7B & 1  & 0  & 1 & 8  & 8 & 1 \\
N & Qwen2.5-7B & 1  & 0  & 0 & 0  & 6 & 0 \\
\midrule
A & Qwen2.5-3B & 7  & 9  & 6 & 1  & 7 & 1 \\
C & Qwen2.5-3B & 2  & 1  & 2 & 7  & 8 & 0 \\
N & Qwen2.5-3B & 1  & 0  & 0 & 0  & 8 & 0 \\
\bottomrule
\end{tabular}
\caption{E1 correct answers per category on our question set (Haiku: mean of three runs; the
standard deviation is 0 except for C ``why'' (0.6) and N, where it is 0.6 to 1.5).}
\label{tab:e1-cat}
\end{table}

\paragraph{What the numbers show.}
\begin{itemize}
  \item \textbf{Numbers are where the workflow wins.} Simple RAG gets at most 6 of 10 changes right,
        because the statement tables are not in the 10-Q passages and the model only finds numbers
        that the text happens to mention. The workflow takes them from XBRL and computes the change
        in SQL, so even Qwen2.5-3B gets 9 of 10. The gap between A and C is larger than the
        confidence intervals for every model.
  \item \textbf{Small models gain the most.} With Qwen2.5-7B the workflow reaches 84\% against 38\%
        for simple RAG on the same passages. A model that runs on a laptop becomes useful once code
        does the numbers.
  \item \textbf{The workflow is stable across runs.} All three Haiku runs of A, A0 and A2 gave exactly
        the same 48 correct answers: the numbers come from SQL, so the model's randomness only
        affects the wording. Simple RAG and closed book vary by one to three answers between runs.
  \item \textbf{``Why'' questions are a draw for Haiku} (11 of 12 for A, 10.3 for C). For Qwen2.5-3B
        they are the weak spot of the workflow (1 of 12): its claims rarely pass the citation check,
        so the system declines instead of guessing.
  \item \textbf{Wrong declines come from small models misreading the question} (wrong company or
        metric): 5--8 of 42 for Qwen on the MacBook and 14 with the Ollama build of Qwen2.5-3B.
        Haiku had none with A.
  \item \textbf{Without retrieval the models mostly cannot answer.} The closed-book system N gets
        31\% with Haiku and 14--18\% with Qwen. Most of its answers say that the model does not know
        the figure, and the few numbers it gets right are hedged guesses (``around \$2 per share'')
        that cite no passage a reader could check. The much larger Opus 5.5 is different: without any
        retrieval it gets 62\% (one run), which we discuss in \cref{sec:contamination}.
  \item \textbf{Version 2 changes nothing here.} Our questions almost never reach document QA, so A2
        has the same accuracy as A, at about 1,100 more input tokens per question.
\end{itemize}

\paragraph{The answer writer (A vs.\ A0).} The writer LLM costs one to two seconds and a few hundred
tokens per answer. It never changed a number that reached the user: Haiku and Qwen2.5-7B passed the
check on every answer (38 of 38 and 36 of 36 in run 1), while Qwen2.5-3B failed it on 25 of 32
answers, mostly because it left out the references. Every rejected answer fell back to the template,
so A and A0 have the same accuracy for all models. The writer makes answers easier to read, but it is
not needed for correctness.

\subsection{E1: accuracy on FinanceBench}

\Cref{tab:e1-fb} shows FinanceBench. Unlike in the preliminary report, every model now answers all
150 questions.

\begin{table}[ht]
\centering\small
\setlength{\tabcolsep}{3.5pt}
\resizebox{\linewidth}{!}{%
\begin{tabular}{llrrrrrrrr}
\toprule
System & Model & Accuracy \% & 95\% CI & Wrong & Declined & Calc. & Analyst & Other & \$ / 1k q \\
\midrule
A workflow    & Haiku 5.5 & 69.1 $\pm$ 1.0 & 63--76 & 22.0 & 24.3 & 35.0 & 29.3 & 39.3 & 0.97 \\
A2 workflow v2 & Haiku 5.5 & \textbf{74.9 $\pm$ 1.0} & 68--81 & 24.3 & 13.3 & 42.3 & 31.3 & 38.7 & 1.19 \\
C simple RAG  & Haiku 5.5 & 54.4 $\pm$ 1.4 & 47--62 & 29.0 & 39.3 & 26.3 & 23.7 & 31.7 & 0.69 \\
N closed book & Haiku 5.5 & 21.3 $\pm$ 1.3 & 15--28 & 61.0 & 57.0 & 4.0  & 19.0 & 9.0  & 0.51 \\
N closed book & Opus 5.5  & 62.7           & 55--70 & 44   & 12   & 41   & 30   & 23   & 17.22 \\
\midrule
A workflow    & Qwen2.5-7B & 16.0 & 11--22 & 36 & 90 & 7  & 5  & 12 & 0 \\
C simple RAG  & Qwen2.5-7B & \textbf{23.3} & 17--31 & 64 & 51 & 10 & 12 & 13 & 0 \\
N closed book & Qwen2.5-7B & 8.0  & 4--13  & 112 & 26 & 0 & 8  & 4  & 0 \\
\midrule
A workflow    & Qwen2.5-3B & 14.7 & 9--21  & 69 & 59 & 4  & 5  & 13 & 0 \\
C simple RAG  & Qwen2.5-3B & \textbf{23.3} & 17--30 & 78 & 37 & 7  & 14 & 14 & 0 \\
N closed book & Qwen2.5-3B & 2.0  & 0--5   & 52 & 95 & 0  & 2  & 1  & 0 \\
\bottomrule
\end{tabular}}
\caption{E1 on FinanceBench (150 questions, 50 per type). Haiku: mean $\pm$ standard deviation over
three runs (counts are means); Opus and Qwen: one run. ``Calc.'' is metrics-generated, ``Analyst''
domain-relevant and ``Other'' novel-generated (correct answers out of 50 each).}
\label{tab:e1-fb}
\end{table}

\begin{itemize}
  \item \textbf{With Haiku the workflow is ahead of simple RAG in every question type,} and version 2
        adds another 5.8 points (69.1 to 74.9\%, with run-to-run deviations of 1.0). The gain is
        almost all on calculations (35.0 to 42.3 of 50), and it comes from fewer declines (24.3 to
        13.3), not from more guessing: the number of wrong answers stays about the same (22.0 and
        24.3). Code still evaluates every formula from cited inputs; only 3 of 900 workflow answers
        used an input that was not printed in a retrieved passage, and all three were flagged.
  \item \textbf{With small models the result flips.} Qwen2.5-7B gets 16\% with the workflow and 23\%
        with simple RAG. Document QA asks the model for structured output: an answer, cited claims
        and, for calculations, each input with its passage and a formula. Qwen often cannot produce
        this from long 10-K passages, so the system declines (90 of 150) instead of guessing. It
        gives far fewer wrong answers than simple RAG (36 against 64), but also fewer right ones.
  \item \textbf{Without retrieval, Haiku and Qwen mostly fail,} and they fail badly: closed-book Haiku
        gives 61 wrong answers, more than any system with retrieval, and Qwen2.5-7B answers almost
        everything from memory and gets 112 of 150 wrong. A small model that guesses from memory is
        the most dangerous setup for financial data.
  \item \textbf{Opus 5.5 without any filing gets 62.7\%,} close to our workflow with Haiku, and 41 of
        50 calculations. We discuss what this means in \cref{sec:contamination}.
\end{itemize}

\subsection{Why version 2 helps: FinanceBench error analysis}
\label{sec:fb-errors}

After the first run we read all 46 answers of version 1 (Haiku, run 1) that were wrong or declined,
together with the gold answer and its evidence page, and assigned each to one cause. Only then did we
build version 2, and we ran it once on all questions without further changes. \Cref{tab:fb-errors}
lists the causes and what version 2 does with the same questions (majority verdict of its three
runs).

\begin{table}[ht]
\centering\small
\begin{tabularx}{\linewidth}{Xrrrr}
\toprule
Cause of the v1 failure & $n$ & v2 correct & v2 wrong & v2 declined \\
\midrule
A financial statement was needed but not retrieved & 14 & \textbf{11} & 3 & 0 \\
Another page was not retrieved (cover page, debt table, earnings release) & 9 & 2 & 0 & 7 \\
Different definition than the gold answer (quick ratio, working capital, capital intensity) & 9 & 1 & 8 & 0 \\
Wrong reasoning or conclusion (including one sign error) & 5 & 2 & 3 & 0 \\
Metric not meaningful for the company (bank, card issuer) & 2 & 1 & 0 & 1 \\
Rounding, or a strict grade (\$8.74 billion for \$8.70 billion) & 2 & 0 & 2 & 0 \\
Wrong line item or scope (segment instead of company) & 2 & 0 & 2 & 0 \\
Incomplete list (acquisitions, regions) & 2 & 1 & 1 & 0 \\
Formula could not be evaluated & 1 & 0 & 0 & 1 \\
\midrule
Total & 46 & 18 & 19 & 9 \\
\bottomrule
\end{tabularx}
\caption{Causes of the 46 wrong or declined answers of version 1 (Haiku, run 1), and the verdict of
version 2 on the same questions (majority verdict of its three runs).}
\label{tab:fb-errors}
\end{table}

\paragraph{What the analysis shows.} Half of the declines (14 of 26) had the same cause: the
calculation needed a line of the balance sheet, income statement or cash-flow statement, and the
search did not return that page. Version 2 fixes 11 of these 14. Over all 150 questions (majority of
three runs), 11 declines and 4 wrong answers of version 1 become correct, while 5 correct answers get
worse, so the net gain is 10 questions. The remaining failures need different fixes:
\begin{itemize}
  \item \textbf{Definitions.} Nine ``wrong'' answers computed a reasonable but different metric, for
        example the quick ratio as (cash + receivables) / current liabilities instead of the gold
        answer's (current assets $-$ inventory) / current liabilities. These are differences of
        convention rather than reading errors; stating the definition in the answer, which the
        workflow does, lets a reader see it.
  \item \textbf{Pages outside the statements.} Questions about the cover page (``Which debt securities
        are registered\ldots''), a debt maturity table or an earnings release still fail because
        search does not find the page; v2 mostly declines them, which is the right behaviour when the
        evidence is missing.
  \item \textbf{Judgement.} Questions such as ``Is 3M a capital-intensive business?'' have no single
        right number. The workflow answers with one metric and the gold answer uses another.
\end{itemize}
Because version 2 was motivated by failures on FinanceBench itself, its FinanceBench score is
somewhat optimistic. We limited this by making one general change, choosing it before running it, and
not tuning anything after the run; the same change makes no difference on our own question set,
where the statements are already covered by SQL.

\subsection{Is FinanceBench in the training data?}
\label{sec:contamination}

Closed-book Haiku gets 21\% of FinanceBench, but closed-book Opus 5.5 gets 62.7\%, including 41 of 50
calculations such as ``What is the FY2018 capital expenditure amount for 3M?'' without seeing the
filing. There are two possible reasons: Opus has memorised the figures of these well-known companies'
10-Ks, or it has seen FinanceBench itself, which has been public since 2023 with its answers.
To tell these apart, we also ran closed-book Opus on our own question set, which did not exist
before this project. It gets 62\% there as well (31 of 50, 95\% interval 48--74), for example 3M's Q2
2023 net sales of \$8,325 million and Tesla's gross profit for Q3 2024, both from memory. It cannot
answer for the latest quarter (Q3 2026), which came after its training data: it lists older quarters
instead and says it does not know whether newer ones exist. So Opus has not simply memorised
FinanceBench. It has memorised the reported figures of large US companies up to its training cut-off,
which our ten companies and FinanceBench's 40 have in common. Its answers are still not usable as they
are: they are marked from memory'', cite no passage, and 12 of the 50 are wrong.

This has two consequences. First, any benchmark built from the filings of large companies before a
model's cut-off, ours included, cannot show on its own whether a large model reads the document or
remembers the answer; a closed-book run of the same model is the baseline that retrieval has to beat.
Second, the value of retrieval for a strong model is less in the score than in what it adds to it:
answers about filings after the cut-off (new quarters, and the thousands of smaller companies in the database),
and a citation to the exact passage or XBRL fact, so that a reader can check every number.

\subsection{Examples}
\label{sec:examples}

\paragraph{Retrieval.} \Cref{tab:retrieval-example} shows the top three passages of each search
method for two ``why'' questions, searched in the MD\&A of the correct 10-Q. For Coca-Cola all three
hybrid results contain the gold reasons (``favorable pricing initiatives'', ``favorable channel and
package mix''); BM25 and semantic search each put a less useful passage first. Searching the whole
filing instead of the MD\&A makes semantic search worse: its top three are the accounting-policy note,
the exhibit index and the equity table, which are about ``Coca-Cola'' and ``Q2 2023'' in general. For
3M no method finds the reasons in the top three: the evidence (``disposable respirator sales
decline'') sits in the segment discussion, while the top results are the general overview, which
uses the question's words (``sales'', ``second quarter'') more often. The workflow then gives the
figures but only part of the explanation, which is what happened in E1.

\begin{table}[ht]
\centering\small
\begin{tabularx}{\linewidth}{llX}
\toprule
Method & Hits & Top three passages (section › heading) \\
\midrule
\multicolumn{3}{l}{\emph{``Why did Coca-Cola's operating income increase in Q2 2023?''}} \\
BM25     & 2 of 3 & Equity income (miss); Operating income: Global Ventures (hit); Operating income: Latin America (hit) \\
Semantic & 2 of 3 & Unit case volume (miss); Operating income: Asia Pacific (hit); Global Ventures (hit) \\
Hybrid   & 3 of 3 & Operating income: Global Ventures, Latin America, Asia Pacific (all hits) \\
\midrule
\multicolumn{3}{l}{\emph{``Why did 3M's sales drop in Q2 2023?''}} \\
BM25     & 0 of 3 & Overview: financial condition; Results of operations: ERP system; Overview \\
Semantic & 0 of 3 & Overview; Overview: financial condition; Financial condition and liquidity \\
Hybrid   & 0 of 3 & Overview: financial condition; Overview; Results of operations \\
\bottomrule
\end{tabularx}
\caption{Real search results for two ``why'' questions (MD\&A only). A hit contains one of the
question's evidence phrases.}
\label{tab:retrieval-example}
\end{table}

\paragraph{Declines.} \Cref{tab:decline-examples} shows how the workflow (Haiku) answered the decline
questions. Each refusal names its reason, which comes from code, not from the LLM: the company is not
in the database, the period would need a 10-K, the company does not report the metric in its XBRL
data, or the passages do not contain the answer. The one ``failure'' is instructive: we expected the
system to decline Coca-Cola's gross margin because margins are not one of our ten metrics, but
document QA found the sentence ``gross profit margin increased to 59.0 percent'' in the 10-Q and
answered correctly with a citation. Our answer key was stricter than necessary here. The closed-book
system, in contrast, answered two decline questions from memory, for example a figure that is only in
a 10-K.

\begin{table}[ht]
\centering\small
\begin{tabularx}{\linewidth}{XX}
\toprule
Question & Answer of the workflow (status) \\
\midrule
What was IKEA's revenue in Q2 2024? & Not enough evidence: no filings for `IKEA' are loaded.
  (abstain) \\
What was Johnson \& Johnson's Q4 2023 revenue? & No filing of Johnson \& Johnson for Q4 2023 is
  loaded (Q4 data is not reported in 10-Q filings). (abstain) \\
What was 3M's R\&D expense in Q2 2023? & 3M does not report research \& development expense in its
  10-Q data (none of the XBRL tags \ldots\ is reported). (abstain) \\
Why did 3M's stock price fall after the earnings release? & The filing passages do not answer this
  question. (abstain) \\
What was Coca-Cola's gross margin percentage in Q2 2023? & Gross profit margin was 59.00 percent
  [C321889]. (answered; expected a decline) \\
\bottomrule
\end{tabularx}
\caption{Answers of the workflow to decline questions of our test set (Haiku 5.5).}
\label{tab:decline-examples}
\end{table}

\subsection{E2: retrieval}

\Cref{tab:e2} compares the three search methods without any LLM. For our ``why'' questions a hit is
a passage that contains one of the evidence phrases; for FinanceBench it is a passage from the
evidence page. MRR is the mean reciprocal rank of the first hit in the top 10.

\begin{table}[ht]
\centering\small
\begin{tabular}{lrrrrrrr}
\toprule
 & \multicolumn{3}{c}{``Why'' questions, MD\&A (12)} & \multicolumn{3}{c}{FinanceBench (150)} & \\
\cmidrule(lr){2-4}\cmidrule(lr){5-7}
Method & hit@1 & hit@5 & MRR & hit@1 & hit@10 & MRR & ms \\
\midrule
BM25            & 0.25 & 0.83 & 0.49 & 0.19 & 0.45 & 0.26 & 9--50 \\
Semantic (MiniLM) & 0.33 & 0.75 & 0.53 & 0.26 & 0.65 & 0.36 & 12--25 \\
Hybrid (RRF)    & \textbf{0.50} & 0.83 & \textbf{0.61} & \textbf{0.29} & 0.65 & \textbf{0.40} & 9--47 \\
\bottomrule
\end{tabular}
\caption{E2 retrieval quality and median search time. Search is deterministic, so there is one run;
with 12 questions one question changes hit@1 by 0.08.}
\label{tab:e2}
\end{table}

Hybrid search puts the right passage first twice as often as BM25 on ``why'' questions and has the
best MRR on both sets, at the same speed, so it is the default. The development questions set the
evidence threshold to a cosine similarity of 0.45: the lowest top score of a development search that
found its evidence was 0.464, and our first guess of 0.55 would have declined 3 of 8 answerable
questions. FinanceBench is harder for search (hit@10 of 0.65) because its evidence is often a table
in a 100-page 10-K, which is exactly the gap that version 2 closes by giving the statements directly.
With only 12 ``why'' questions the differences between the methods are small in absolute terms (one
question is 0.08), so we read E2 as ``hybrid is at least as good as either method alone'' rather than
as a precise ranking.

\subsection{E3: latency per step}

\begin{table}[ht]
\centering\small
\begin{tabular}{lrr}
\toprule
Step & Number question (ms) & ``Why'' question (ms) \\
\midrule
LLM parse & 4,301 & 4,341 \\
SQL: company / filing / figures & 1.3 / 0.7 / 1.1 & 1.7 / 0.6 / 2.9 \\
Question embedding / vector search / ranking & -- & 2.9 / 3.0 / 0.0 \\
LLM claims & -- & 15,439 \\
LLM writer & 2,940 & 5,838 \\
\midrule
Total (median) & 7,072 & 25,726 \\
\bottomrule
\end{tabular}
\caption{E3 median time per step on the c7i.8xlarge, with Qwen2.5-3B on the CPU (Ollama), over the
answered questions of one run (18 number and 11 ``why'' questions).}
\label{tab:e3}
\end{table}

Our own steps (SQL and search) take less than 10 ms per question. Everything else is LLM time: an LLM
call takes 100 to 1,000 times longer than the rest of the pipeline. Running Qwen2.5-3B on the CPU of
the instance works without any API key, but it is slower and less accurate (54\%) than the same
model on a laptop's Apple GPU (62\%, 9 s per question including the SSH tunnel).

\subsection{E4: load test}
\label{sec:e4}

\paragraph{The pipeline in one process.} With the LLM steps replaced by fixed outputs, only SQL,
search and code run (\cref{tab:e4}).

\begin{table}[ht]
\centering\small
\begin{tabular}{rrrr}
\toprule
Concurrent users & Questions / s & Median ms & p95 ms \\
\midrule
1  & 82.8  & 11  & 17 \\
10 & 284.5 & 34  & 46 \\
50 & 221.5 & 206 & 322 \\
\bottomrule
\end{tabular}
\caption{E4 throughput and latency of one Python process on the c7i.8xlarge, without the LLM
(30 seconds per step).}
\label{tab:e4}
\end{table}

One Python process serves about 280 questions per second at 10 users. At 50 users it is saturated
and latency grows, because one Python process runs Python code on one core at a time; more worker
processes, or several instances behind a load balancer, are the next step.

\paragraph{The HTTP API on three instance sizes.} The second part of E4 tests the deployed API with
up to 1,000 users on c7i.large, c7i.xlarge and c7i.2xlarge (\cref{sec:experiments-e4}). Two changes
were needed before it could work at all. A request used to open its database connection before the
LLM parse step, so 1,000 requests waiting for an LLM would have held 1,000 connections, ten times
Postgres' default limit of 100; the connection is now opened after the parse, and the limit is
raised to 500. And FastAPI serves synchronous requests from a pool of 40 threads per process, which
would cap the number of requests waiting for an LLM at the same time; the test server raises it.

\resultsplaceholder{E4 HTTP results: answers per second, median / p95 / p99 latency and error rate
for 1--1,000 users on c7i.large, c7i.xlarge and c7i.2xlarge, with no LLM wait and with a 5-second
LLM wait. To be run once the AWS vCPU quota allows the three extra instances.}

\subsection{E5: Spark scaling}
\label{sec:e5}

\begin{figure}[ht]
\centering
\begin{tikzpicture}
\begin{axis}[
  width=0.62\linewidth, height=55mm, xmode=log, log basis x=2,
  xtick={1,2,4,8,16,32}, xticklabels={1,2,4,8,16,32},
  xlabel={Spark cores}, ylabel={Speed-up}, ymin=0, ymax=12,
  legend pos=north west, legend style={font=\small}, grid=major, grid style={gray!20}]
\addplot[wf,thick,mark=*] coordinates {(1,1.0) (2,1.9) (4,3.2) (8,5.2) (16,7.7) (32,8.5)};
\addplot[cb,thick,mark=square*] coordinates {(1,1.0) (2,1.8) (4,3.2) (8,5.2) (16,7.0) (32,7.0)};
\addplot[gray,dashed,domain=1:12,samples=2] {x};
\legend{Text job, Facts job, Ideal (linear)}
\end{axis}
\end{tikzpicture}
\hspace{4mm}
\begin{minipage}[b]{0.3\linewidth}\small
\begin{tabular}{rrr}
\toprule
Cores & Text s & Facts s \\
\midrule
1 & 170 & 284.8 \\
2 & 90 & 157.3 \\
4 & 53 & 88.3 \\
8 & 33 & 55.0 \\
16 & 22 & 40.6 \\
32 & 20 & 40.6 \\
\bottomrule
\end{tabular}
\vspace{6mm}
\end{minipage}
\caption{E5 speed-up of the two Spark jobs on one c7i.8xlarge (32 vCPUs = 16 physical cores with
two threads each). The text job ran on 80 filings (2,458 passages), the facts job on all 20,609
company files; one run per core count.}
\label{fig:e5}
\end{figure}

Both jobs scale well up to 8 cores (5.2 times faster) and still gain up to 16. Going from 16 to 32
adds little, for two reasons: the 32 vCPUs are only 16 physical cores, and Spark's fixed start-up
time weighs more on these small test jobs. On the full batches of 2,000 filings the instance reached
130 to 190 passages per second.

\subsection{E6: updates with new filings}
\label{sec:e6}

\Cref{tab:e6} shows the update run of 26 November 2026. Since the last load (filings up to 20
November) the SEC had received 92 new 10-Qs, including three companies that were not in the database
yet. The whole update took five minutes. 53 of the 92 filings had an MD\&A section that our splitter
recognised; they gave 1,950 new passages. (The share is the same as in the full run, where 72\% of
the 10-Qs got text.)

\begin{table}[ht]
\centering\small
\begin{tabularx}{\linewidth}{Xr}
\toprule
Step (c7i.8xlarge, 32 vCPUs) & Seconds \\
\midrule
1. Download the SEC bulk files again (\code{companyfacts.zip}, \code{submissions.zip}, 2.8 GB) & 114 \\
2. Spark facts job on all 20,609 companies (38 s inside Spark; the rest is starting the container and Spark) & 124 \\
3. Load facts: 92 new filings, 3 new companies, 2,744 new figures & 17 \\
4. Manifest of the new filings & 7 \\
5. Download the 92 10-Qs & 13 \\
6. Spark text job: split and embed (1,950 passages) & 25 \\
7. Load the passages & 2 \\
\midrule
Total & 301 \\
\bottomrule
\end{tabularx}
\caption{E6: one update run, adding the 10-Qs filed from 21 to 26 November 2026.}
\label{tab:e6}
\end{table}

\paragraph{How much does a busy day need?} Late November is a quiet time. \Cref{fig:filings-per-day}
shows how many 10-Qs the SEC receives per day: about 17,500 a year, but very unevenly. On a median
day there are only 16, while in the two weeks after each quarter's deadline there are several hundred
per day, with a maximum of 718 (7 November 2024). For the busiest day we can estimate the update time
from the measured rates. The download is limited by the SEC to under 10 requests per second, so 718
filings take about 90 seconds on any instance. At about 37 passages per filing, they give about
26,600 passages to embed: about 3 minutes at the 130--190 passages per second of the c7i.8xlarge,
or about 16 minutes on a c7i.large with 2 vCPUs (27 passages per second in E5). Adding the bulk
download and the facts job (2 to 3 minutes), even the busiest day of the year is updated in under
10 minutes on the large instance and in about 25 minutes on the smallest one. A daily update
therefore does not need the large instance at all: a small instance that runs once a night is enough,
and the large one is only needed for the first full load, which processed the 140,000 10-Qs of seven
years in about 9 hours.

\begin{figure}[ht]
\centering
\begin{tikzpicture}
\begin{axis}[
  width=0.95\linewidth, height=45mm, ybar interval, ymode=log, log origin=infty, ymin=0.8, ymax=1000,
  xlabel={10-Qs filed on one day}, ylabel={Days (2024--2026)},
  xmin=0, xmax=800, enlarge x limits=false, x tick label as interval=false,
  xtick={0,100,200,300,400,500,600,700,800}, xticklabel style={font=\footnotesize},
  yticklabel style={font=\footnotesize}, grid=major, grid style={gray!20}]
\addplot[fill=wf,draw=white] table[x=bin,y=days] {figures/filings_hist.dat};
\end{axis}
\end{tikzpicture}
\caption{E6: how many 10-Qs the SEC receives per day, from \code{submissions.zip} (2024 to November
2026, 720 working days). Most days have under 50; about 20 days a year have over 300.}
\label{fig:filings-per-day}
\end{figure}
